\documentclass[11pt]{article}
\usepackage[margin=0.9in]{geometry}
\usepackage{amsmath,amssymb,amsthm,booktabs,mathrsfs}
\usepackage[T1]{fontenc}
\usepackage{lmodern,microtype}
\usepackage[colorlinks=true,linkcolor=blue,urlcolor=blue]{hyperref}
\newtheorem{theorem}{Theorem}[section]
\newtheorem{lemma}[theorem]{Lemma}
\newtheorem{proposition}[theorem]{Proposition}
\newtheorem{corollary}[theorem]{Corollary}
\theoremstyle{remark}
\newtheorem{remark}[theorem]{Remark}
\newcommand{\norm}[1]{\left\|#1\right\|}
\title{Toward a four-dimensional Yang--Mills construction\\
\large Notebook 33: gauge-invariant determinant densities and Cartan probes}
\author{Research working document}
\date{17 October 2026}
\begin{document}
\maketitle
\begin{abstract}
We identify a constant-curvature calculation with the actual
Maxwell coefficient of the leading quantum correction.
The starting point is the exact vector, based-ghost and
effective-Hessian determinant identity of Notebook 32.

At a fixed composite endpoint, the matrices transform by
background-gauge conjugation. Their diagonal color traces
therefore define individually gauge-invariant analytic
densities. Fixed-endpoint gaps and exponentially localized
source kernels make their matrix logarithms local, uniformly
in compatible volume. A rooted-density Ward argument proves
that their quadratic affine coefficient is precisely the
second-displacement moment of the total Hessian. This
resolves the ambiguity that would arise from evaluating
an arbitrary integration-by-parts density on an unbounded
affine connection.

Commuting affine connections are exactly Wilson stationary
and are preserved by the actual centered ordinary averaging
maps. An exponential-kernel uniqueness argument identifies
their polynomial stationary jets with the constructed
minimum section. The envelope identity extends the
zero-current conclusion to the effective classical actions.
The constraint multiplier consequently vanishes through
the background orders needed for the Maxwell calculation.

The resulting coefficient is an explicit diagonal trace
expression on a constant Cartan-curvature probe, using
the actual composite constraint and its normalized
determinants. For the bare composite it involves the
Wilson Hessian directly. The coarse contribution already
has a uniform bound. We compute a specified massive
continuum proper-time reference, including its positive
logarithmic coefficient and a uniform finite-part bound.
An exact heat-integral formula for the actual lattice
coefficient isolates sufficient estimates on three
heat-time intervals. Those estimates for the actual
family remain incomplete. A closed-walk expansion proves
the short-time bound for the actual based-ghost part,
uniformly in endpoint and preceding depth. The vector
part, the other heat-time intervals and the subsequent
full-measure and continuum arguments remain.
\end{abstract}

All new arguments below are working proofs checked against
the stated inputs from Notebooks 15--32. The matrix-logarithm
locality estimates fix the endpoint before taking a volume
limit. They are not estimates uniform in the ultraviolet
depth. An affine field is used only through polynomial
extensions of localized derivatives; no periodic family
of continuously varying magnetic flux is assumed.


% No preamble. Prefix gdd-. Fixed endpoint, actual determinant densities.
\section{Gauge-invariant diagonal determinant densities}
\label{gdd-section}

Fix a composite endpoint \(K=\ell^k\), a preceding homogeneous
depth \(M\), and the actual finite periodic construction. The
endpoint is fixed throughout this section. Constants may depend
on \(K\); they will be uniform in the permitted volume and in
the preceding depth where the upstream classical bounds are.
Choose the fixed sufficiently large scalar \(a>0\) from the
coarse determinant estimate of Notebook 32.

At the actual stationary field \(U_*(W)\), use its left-relative
objects
\[
 C=f'(0),\quad s'=C^T\lambda,\quad
 \mathcal L=s''-\lambda\cdot f'',\quad
 G_b=D_b^TD_b,\quad
 Q_b=\mathcal L+D_bD_b^T+aC^TC,\quad
 P_c=q_{\mathrm{eff}}+aI.
\]
The exact finite identity already established is
\begin{equation}
 D_{M,K}(W)=\mathcal N_0\left[
       \log\det G_b-\frac12\log\det Q_b
                           +\frac12\log\det P_c\right].
 \label{gdd-exact-total}
\end{equation}
Here \(\mathcal N_0\) subtracts the value at \(W=1\).
The matrices act, respectively, on noncenter scalar sites,
current bonds, and output bonds, including their color coordinates.

\subsection{Background covariance of the actual matrices}
\label{gdd-covariance-section}

For a fixed current-site gauge \(u\), let \(O_E\) denote
the block-diagonal adjoint action at each current bond's
initial site, \(O_S\) the corresponding action at noncenter
scalar sites, and \(O_C\) its action at output bond initial
centers. These matrices are orthogonal for real gauges.
They are complex orthogonal on the holomorphic continuation.
The transformation of a left-relative variation is linear:
\[
 (e^{ix}U)^u=e^{iO_Ex}U^u.
\]
The actual ordinary average is equivariant, so its relative
target logarithm transforms by \(O_C\). Local gauge invariance
of the two actual classical actions consequently gives
\begin{align}
 C(U^u)&=O_C C(U)O_E^T,&
 \lambda(W^{u|_{\rm centers}})&=O_C\lambda(W),\notag\\
 \mathcal L(U^u)&=O_E\mathcal L(U)O_E^T,&
 D_b(U^u)&=O_E D_b(U)O_S^T. \label{gdd-transformations}
\end{align}
For the Lagrangian identity, differentiate
\[
 s_{U^u}(O_Ex)=s_U(x),\qquad
 f_{U^u}(O_Ex)=O_Cf_U(x).
\]
The second equality and the multiplier transformation retain
the entire \(\lambda\cdot f''\) term. No unconstrained
stationarity is assumed here.

The based scalar domain is unchanged by these transformations:
zero values at all \(K\)-centers remain zero. The effective
relative Hessian transforms by \(O_C\), since its target
relative variable also transforms linearly. Thus
\begin{equation}
 G_b(U^u)=O_SG_b(U)O_S^T,\quad
 Q_b(U^u)=O_EQ_b(U)O_E^T,\quad
 P_c(W^{u|_{\rm centers}})=O_CP_c(W)O_C^T.
 \label{gdd-matrix-covariance}
\end{equation}
These identities concern the fixed scalar completion \(aI\).
They do not assert such covariance for an arbitrary
background-dependent completion.

\subsection{Fixed-endpoint locality of the matrix logarithms}
\label{gdd-log-locality-section}

We need locality of a diagonal logarithm, not merely locality
of the trace of its derivatives. The required fixed-endpoint
spectral and kernel inputs follow from the preceding modules,
as follows.

At zero background, \(G_b\) has a positive gap independent
of volume when \(K\) is fixed. Indeed, join every site to
its block center by a path of length at most \(c_dK\).
For a scalar field vanishing at centers, Cauchy--Schwarz
along these paths and summation in a block give
\begin{equation}
 \|\phi\|_2^2\le c_d K^{d+1}\|D_b\phi\|_2^2.
 \label{gdd-pinned-Poincare}
\end{equation}
The crude exponent is sufficient here. This is compatible
with the soft-mode upper bound in Notebook 32 and supplies
no gap uniform as \(K\) grows.

The flat constrained contour precision has its previously
proved positive gap at fixed \(K\), uniformly in preceding
depth and volume. On \(\ker C\), decompose a vector into
its component in \(\operatorname{ran}D_b\) and its orthogonal
complement. The form \(\mathcal L\) annihilates the first.
For a horizontal vector \(v\in\ker C\), replace it by
\(v+D_bTv\) to make its contour derivative zero.
Its \(\mathcal L\) value is unchanged, and the orthogonal
projection of this replacement back to the horizontal
space is \(v\). The contour precision bound therefore
controls \(\|v\|_2^2\). Its passage from retained
coordinates to current bonds uses the fixed-\(K\)
bounded insertion \(K_0\); this only changes the constant.
On vertical vectors \(D_b\phi\), the divergence penalty
has the lower bound
\[
 \|D_b^TD_b\phi\|_2^2
       \ge \lambda_{\min}(G_b)\|D_b\phi\|_2^2.
\]
Thus \(\mathcal L+D_bD_b^T\) has a uniform fixed-\(K\)
gap on \(\ker C\).

The actual adjusted minimum-section derivative
\[
 H_b=H-D_bG_b^{-1}D_b^TH
\]
has \(CH_b=I\), \(D_b^TH_b=0\), and a bounded operator norm
at fixed \(K\). The latter follows from the established
minimum-response bound and the orthogonal projection
in this formula. Write \(x=\xi+H_by\), \(C\xi=0\).
The exact effective Schur identity gives
\[
 x^TQ_bx
  =\xi^T(\mathcal L+D_bD_b^T)\xi
                         +y^T(q_{\mathrm{eff}}+aI)y.
\]
Since \(q_{\mathrm{eff}}+aI\) has a fixed positive gap,
this proves a flat gap for \(Q_b\), uniform in volume
and preceding depth at fixed \(K\). Their upper spectral
bounds and exponentially weighted matrix bounds follow
from the actual strict classical families, the finite-range
constraint at fixed \(K\), and the local gauge generator.
\(P_c\) already has the stronger uniform bounds of Notebook 32.

Compose these matrices with the actual stationary section.
Its strict support kernels, together with those of the
classical actions and local charts, give source-rooted
matrix derivative bounds through order two. A sufficiently
small common complex output ball preserves the flat
gaps in operator norm. This ball may depend on \(K\),
but not on volume or preceding depth.

Here is an elementary logarithm argument making the
locality consequence precise. For any one of these
matrix families \(P(w)\), choose a fixed scalar \(s\)
larger than its flat upper spectral bound and put
\[
 Z(w)=I-P(w)/s.
\]
At zero \(\|Z(0)\|_{2\to2}<1\), uniformly over the
volumes under consideration. The exponentially weighted
kernel bound implies that conjugating \(Z\) by a small
exponential distance weight changes its operator norm
by an arbitrarily small amount. First use bounded
truncations of the distance weight; their uniform
Lipschitz bound gives the same estimate, and then pass
to kernel bounds. No unbounded similarity is assumed
to act on all of \(\ell^2\). Shrink the background
ball as well. The conjugated operators then have norm
at most one fixed \(r_K<1\). Therefore
\begin{equation}
 \log P(w)=(\log s)I-\sum_{n\ge1}\frac{Z(w)^n}{n}
 \label{gdd-log-series}
\end{equation}
converges with these exponential weights.
Termwise differentiation through order two is permitted:
there are at most \(n\) or \(n^2\) insertions, absorbed
by the geometric factor \(r_K^{n-2}\).
The actual source-rooted bounds attach every background
source to the matrix-product chain. Exponential weights
then connect the diagonal matrix coordinate and all
sources. Conversion from weighted \(\ell^2\) to row
sums spends a smaller exponent and a finite geometric
sum. The internal fiber per output block is finite
at fixed \(K\).

In particular, if a diagonal color trace is assigned
to a fixed output block \(x\), its quadratic kernel
\(k_x(b,c)\) obeys, for some \(\kappa_K>0\),
\begin{equation}
 \sum_{b,c}e^{\kappa_K\tau(x,b,c)}|k_x(b,c)|\le C_K.
 \label{gdd-rooted-kernel}
\end{equation}
Here the sum includes directions and colors, \(\tau\)
is a connecting-tree distance in output units, and all
finite multiplicities of current coordinates are included.
Periodic bounds use shortest lifted distances. No extensive
logarithm is taken to an infinite-volume limit.

\subsection{A density with its own gauge identity}
\label{gdd-density-section}

Assign every noncenter current site to its unique centered
\(K\)-block. Assign each positive current bond to the block
containing its initial site, and each output bond to its
initial output site. Define
\begin{align}
 \rho_x(w)=\mathcal N_0\bigg[
 &\sum_{\substack{z\ {\rm noncenter}\\z\ {\rm in\ block}\ x}}
                   \operatorname{tr}_{\mathfrak g}(\log G_b)_{zz}
 -\frac12\sum_{\substack{e\ {\rm current}\\e_-\ {\rm in\ block}\ x}}
                   \operatorname{tr}_{\mathfrak g}(\log Q_b)_{ee}
 \notag\\
 &+\frac12\sum_{\substack{c\ {\rm output}\\c_-=x}}
                   \operatorname{tr}_{\mathfrak g}(\log P_c)_{cc}
 \bigg].
 \label{gdd-density}
\end{align}
All matrices in this formula are evaluated at \(U_*(e^{iw})\)
or its output \(e^{iw}\), as appropriate. On the finite
regulator,
\begin{equation}
                         D_{M,K}(e^{iw})=\sum_x\rho_x(w)
 \label{gdd-density-total}
\end{equation}
exactly.

Functional calculus respects the similarities in
\eqref{gdd-matrix-covariance}. A diagonal color trace is
unchanged by its site's orthogonal conjugation. Moreover,
the actual minimum at gauge-transformed output is gauge
equivalent to the transformed old minimum, on the established
local germ. To see this, transform the constraint and action,
then use the proved uniqueness of the local constrained
critical orbit and its selected contour representative.
The comparison need only hold on a finite real germ;
analyticity supplies the local identities used below.
It follows that \emph{each} \(\rho_x\), and each of its
three displayed groups of terms, has its own local
output gauge identity. This conclusion uses more than
gauge invariance of their sum.

The block assignment is translation covariant. Local
convergence of the matrices and their derivatives, together
with \eqref{gdd-log-series} and its uniform geometric tail,
passes these densities and their quadratic kernels to
any jointly compatible infinite system at fixed endpoint.
There they are translation covariant, satisfy
\eqref{gdd-rooted-kernel}, and retain the local gauge identity.
No uniqueness of different structural limits is needed.

\subsection{An affine-density formula for the total Maxwell moment}
\label{gdd-moment-section}

\begin{lemma}[Gauge-invariant rooted density and total moment]
\label{gdd-moment-lemma}
Let \(\rho_x\) be translation-covariant analytic scalar
densities on an infinite lattice, with exponentially
rooted quadratic kernels as in \eqref{gdd-rooted-kernel}.
Assume every \(\rho_x\) is invariant under local gauge
changes on its germ, and under global semisimple color
conjugation. Fix one color and distinct directions
\(\mu,\nu\). For the formal affine component
\[
                  A_\mu(z)=t z_\nu,\qquad A_\sigma=0
                    \quad(\sigma\ne\mu),
\]
the quadratic density coefficient is
\begin{equation}
 [t^2]\rho_0(A)
       =-\frac14\sum_r r_\nu^2 K_{\mu\mu}(0,r),
 \label{gdd-affine-moment}
\end{equation}
where \(K=\sum_x D^2\rho_x(0)\) is the total Hessian
kernel in that color. Both sides are absolutely convergent
kernel expressions. The left side does not mean evaluation
of an analytic map at an unbounded nonzero field.
\end{lemma}

\begin{proof}
Global semisimple color invariance kills each density's
linear derivative. Differentiating its local gauge identity
at the origin therefore shows that its Hessian annihilates
each compactly supported infinitesimal gradient in either
slot. The exponential kernel bound extends this identity
to polynomial gauge parameters by absolute summation.
In particular it annihilates a constant one-form in any
one color and direction: use the linear scalar gauge
parameter with that gradient.

Write \(k_0(b,c)\) for the selected same-component Hessian.
Its row and column sums are zero. Consequently
\[
 \frac12\sum_{b,c} k_0(b,c)b_\nu c_\nu
   =-\frac14\sum_{b,c}k_0(b,c)(c_\nu-b_\nu)^2.
\]
Every interchange is allowed by exponential rooted
summability. Translation covariance gives
\[
 K(0,r)=\sum_x k_x(0,r)=\sum_x k_0(-x,r-x).
\]
The change of indices \(b=-x,\ c=r-x\) turns the last
second moment into the right side of
\eqref{gdd-affine-moment}. The same calculation with
the translated affine field shows that the coefficient
is independent of the selected density root.
\end{proof}

For the actual correction, the established Maxwell
normalization is
\(\beta(D)=-\tfrac12\sum_r r_\nu^2K_{\mu\mu}(0,r)\),
in one normalized color. Thus the constructed density gives
\begin{equation}
                  \boxed{\ \beta(D_{M,K})
                            =2[t^2]\rho_0(tz_\nu e_\mu)\ }.
 \label{gdd-beta-density}
\end{equation}
The density construction applies separately to the
three groups of terms in \eqref{gdd-density}, with the
corresponding second-moment extraction on each.
When the total has the established spatial symmetries,
this is its Maxwell coefficient.

The individual gauge identity is essential. For example,
the Maxwell action may be written either as its curvature
square density or, after summation by parts, as a density
of the form \(\tfrac12 A\cdot d^*dA\). On an affine
constant-curvature field the latter vanishes pointwise,
while the former does not. Their totals agree on compact
fields. The diagonal densities above avoid this ambiguity
because their own gauge identities imply the row and
column cancellations used in the proof.

\subsection{Scope}
\label{gdd-scope}

The result identifies a specific, locally gauge-invariant
trace density with the total Maxwell moment. It does not
estimate that moment uniformly as \(K\) grows. The gap,
analytic radius and decay exponent used for \(G_b,Q_b\)
are permitted to deteriorate with \(K\). No positive
coefficient or comparison with a massive reference
follows from the density identity alone.


% No preamble. Prefix cdmom-. Polynomial jets, not a periodic flux family.
\section{Constant Cartan curvature and the actual stationary jets}
\label{cdmom-section}

We now use the diagonal-density construction to make a
constant-curvature calculation an exact test of the actual
Maxwell coefficient. All substitutions at an unbounded
affine field mean substitutions into exponentially localized
Taylor kernels at zero, through the required finite order.

\subsection{Exact local Wilson and averaging identities}
\label{cdmom-affine-section}

Choose commuting traceless Hermitian matrices \(c_\mu,B_{\mu\nu}\)
in a Cartan subalgebra of \(\mathfrak{su}(N)\), with
\(B_{\mu\mu}=0\), and put
\begin{equation}
 A_\mu(x)=c_\mu+\sum_\nu B_{\mu\nu}x_\nu,\qquad
 U_\mu(x)=e^{iA_\mu(x)}.
 \label{cdmom-affine}
\end{equation}
The plaquette in the \(\mu,\nu\) plane is the constant matrix
\[
 e^{iF_{\mu\nu}},\qquad F_{\mu\nu}=B_{\nu\mu}-B_{\mu\nu}.
\]
The full Wilson Euler derivative is zero. For each varied
bond and transverse direction, the two incident plaquettes
give opposite transported sine terms. They agree after
transport because all background links and plaquettes
commute. The cancellation occurs as a Lie-algebra-valued
identity, including against variations outside the Cartan
subalgebra. This is a local infinite-lattice identity;
an artificial finite-box boundary is not included.

For a centered primitive block, write \(x=y+r\).
Every coordinate-order root contour has transporter
\[
 \exp\{i[r\cdot A(y)+q_\Gamma(B,r)]\},
\]
where \(q_\Gamma\) is independent of \(y\).
The selected mean is endpoint equivariant and conjugation
equivariant. Its value on a tuple of diagonal matrices
is diagonal: conjugation by every diagonal group element
fixes the tuple, hence its output. Therefore the actual
root transporter has the form
\begin{equation}
 P_\ell(U;y,y+r)=e^{ir\cdot A(y)}p_\ell(B,r).
 \label{cdmom-contour}
\end{equation}
No log-arithmetic formula for the selected mean is assumed.

The \(p_\ell\) factors cancel in the ordinary path product
between neighboring blocks. Since \(B_{\mu\mu}=0\), this
product is
\[
 U_x=e^{i\ell A_\mu(y)}
       \exp\!\left\{i\ell\sum_\nu
                 (B_{\mu\nu}-B_{\nu\mu})r_\nu\right\}.
\]
Its reference center string is \(e^{i\ell A_\mu(y)}\).
The defining ordinary average takes the centered average
of the relative logarithms, and \(\sum_r r_\nu=0\).
It follows that the actual primitive and iterated maps obey
\begin{equation}
 M_\ell(U)(y,y+\ell e_\mu)=e^{i\ell A_\mu(y)},\qquad
 F_K(U)(y,y+Ke_\mu)=e^{iK A_\mu(y)},\quad K=\ell^k.
 \label{cdmom-average}
\end{equation}
Thus the output affine curvature is \(K^2F\).
These are identities of local analytic germs, hence
identities of formal curvature jets at each lattice site.
They do not require a common bounded logarithm of the
affine field on the whole infinite lattice.

Apply the actual direct-\(K\) root gauge
\(u(x)=P_K(U;y,x)\) in the block centered at \(y\).
It is the identity at centers. Endpoint covariance makes
the transformed contour equal to the identity and leaves
the composite output unchanged. Wilson stationarity
is preserved by gauge invariance. At every fixed order
in a curvature parameter, the rooted candidate has
polynomial growth in the block center. The within-block
coordinate set is finite at fixed \(K\).

\subsection{Why the formal candidate is the actual minimum jet}
\label{cdmom-jet-section}

\begin{lemma}[Polynomial extension of the stationary recursion]
\label{cdmom-polynomial-lemma}
Suppose an actual stationary equation in retained coordinates is
\[
 \mathcal A z+\mathcal B w+\mathcal N(w,z)=0,\qquad
                     \mathcal N=O((w,z)^2).
\]
Assume \(\mathcal A,\mathcal A^{-1},\mathcal B\) have
exponentially summable kernels, the inverse identities
hold on finitely supported fields, and the mixed Taylor
kernels of \(\mathcal N\) through a specified finite
order have exponential rooted-tree bounds. Then its
analytic solution jets extend uniquely to polynomial-growth
input jets. Any polynomial-growth formal solution has
those same coefficients through that order.
\end{lemma}

\begin{proof}
For a polynomial-growth sequence \(v\),
\[
 (1+|j|)^m\le(1+|i|)^m(1+d(i,j))^m
\]
shows that convolution by an exponentially summable
kernel is absolutely defined and preserves polynomial
growth. Exponential summability also permits the two
convolutions in the inverse identities to be interchanged.
Those identities therefore hold on polynomial-growth
sequences. In particular \(\mathcal A\) has no nonzero
polynomial-growth nullvector.

For a mixed kernel of any fixed finite degree, bound
each input polynomial by a polynomial in the output
coordinate times a polynomial in the joining-tree lengths.
The exponential tree bound absorbs the latter factors.
Thus the nonlinear Taylor coefficients are well defined.
At order \(n\) the stationary equation is
\[
 \mathcal A z_n=h_n(w_1,\ldots,w_n,z_1,\ldots,z_{n-1}),
\]
and has the unique polynomial-growth solution
\(\mathcal A^{-1}h_n\). These are also the equations
obtained by differentiating the analytic bounded-field
solution. Induction proves the assertion.
\end{proof}

For the actual constrained stationary branch at a fixed
endpoint, these are established inputs: the flat retained
precision has its exponential inverse, the affine
constraint chart is local, and the classical scalar
and stationary support families have exponentially
localized mixed derivatives. The finite fiber dimension
is per block; the whole lattice is not finite-dimensional
in this application. Only finite-order jets are being
extended.

For the bare Wilson action, the rooted affine candidate
above solves the constraint, contour and stationary
equations coefficientwise, with unconstrained current
and multiplier equal to zero. Lemma~\ref{cdmom-polynomial-lemma}
therefore identifies it with the actual stationary
section's polynomial jets. In particular this applies
through the two background orders needed below.

\subsection{The effective classical action is also stationary on this family}
\label{cdmom-effective-section}

The effective-action interface can be checked using the
actual envelope identity, without replacing its Hessian
by a bare Wilson Hessian.

Fix any permitted finite homogeneous depth \(M\).
The action \(S_M\) is the actual local minimum value of
the bare Wilson action under \(F_M\). Its differential
in left-relative output coordinates equals the target
multiplier of that minimum. This identity is an equality
of the established analytic kernels. It extends to
polynomial input jets by their exponential bounds.

Apply the preceding bare-action argument at the fixed
endpoint \(M\), with fine curvature equal to the prescribed
current curvature divided by \(M^2\). The explicit fine
candidate is Wilson stationary, so its multiplier jets
are zero. Consequently
\begin{equation}
          D_{\rm left}S_M(U^{\rm affine}(t))=0
                 \quad\hbox{as formal jets through order two}.
 \label{cdmom-effective-stationarity}
\end{equation}
The identity holds in every current-link component;
it is not just a derivative along the one-parameter family.
Indeed the envelope identifies the whole target multiplier,
and the unconstrained fine Euler derivative is zero.

At the next fixed endpoint \(K\), take the current affine
field with curvature divided by \(K^2\), and root it by
the actual direct-\(K\) based gauge. The classical local
gauge identity, also extended coefficientwise by exponential
bounds, preserves \eqref{cdmom-effective-stationarity}.
This candidate satisfies the actual \(S_M\) constrained
stationary equation with zero multiplier through the same
orders. Apply Lemma~\ref{cdmom-polynomial-lemma} to the
retained precision for \(S_M\). The candidate has exactly
the actual current stationary jets, and
\begin{equation}
             \lambda_{M,K}(t)=O(t^3)
 \label{cdmom-multiplier-jets}
\end{equation}
in this formal sense.

All assertions so far fix each endpoint when invoking
its inverse. They do not infer a gap uniform in \(M\)
or \(K\) for a one-shot fine operator. However, the
current-coordinate classical and stationary Taylor
kernels at fixed \(K\) have the already established
uniform bounds in preceding depth. Their polynomial
moments consequently pass to any varying-depth
structural limit by exponential domination.
Thus \eqref{cdmom-effective-stationarity} and
\eqref{cdmom-multiplier-jets} also hold in those joint
systems. The effective Hessian \(s''\) still depends
on \(S_M\); this result does not make it the bare
Wilson Hessian.

\subsection{An exact constant-curvature expression for the Maxwell coefficient}
\label{cdmom-determinant-section}

Fix a trace-orthonormal Cartan color \(H\), and distinct
directions \(\mu,\nu\). The output probe is
\[
 w_\mu(n;t)=t n_\nu H,\qquad w_\sigma=0\quad(\sigma\ne\mu).
\]
Its unrooted current affine representative is
\begin{equation}
 U_\mu^{\rm aff}(x;t)=\exp(i t x_\nu H/K^2),\qquad
 U_\sigma^{\rm aff}=I\quad(\sigma\ne\mu).
 \label{cdmom-probe}
\end{equation}
By \eqref{cdmom-average}, its output is exactly the
chosen formal probe. Its oriented \(\mu,\nu\) plaquette
has angle \(-tH/K^2\).

On this family write
\begin{align}
 G(t)&=D_b(U^{\rm aff}(t))^TD_b(U^{\rm aff}(t)),\notag\\
 V_a(t)&=s''(U^{\rm aff}(t))
           +D_b(U^{\rm aff}(t))D_b(U^{\rm aff}(t))^T
           +aC(U^{\rm aff}(t))^TC(U^{\rm aff}(t)),\notag\\
 P(t)&=q_{\rm eff}(w(t))+aI.
 \label{cdmom-probe-matrices}
\end{align}
All are formal matrix jets around their flat positive
operators. In particular no invertibility or positivity
on a nonzero global magnetic background is assumed.
Equation \eqref{cdmom-multiplier-jets} removes the
\(-\lambda\cdot f''\) term through second order.
The remaining Hessian \(s''\) retains the actual action.

Use the same block-zero site and bond assignments as
in the diagonal-density construction. Its pointwise
gauge invariance allows removal of the root gauge,
coefficientwise: the local gauge identities extend to
these polynomial gauge jets by the same exponential
summation. The density--moment lemma therefore gives
\begin{equation}
 \boxed{\begin{split}
 \beta(D_{M,K})=[t^2]\bigg\{
 &2\sum_{\substack{z\ {\rm noncenter}\\z\ {\rm in\ block}\ 0}}
                   \operatorname{tr}_{\mathfrak g}(\log G(t))_{zz}\\
 &-\sum_{\substack{e\ {\rm current}\\e_-\ {\rm in\ block}\ 0}}
                   \operatorname{tr}_{\mathfrak g}(\log V_a(t))_{ee}\\
 &+\sum_{\substack{c\ {\rm output}\\c_-=0}}
                   \operatorname{tr}_{\mathfrak g}(\log P(t))_{cc}
                      \bigg\}.
 \end{split}}
 \label{cdmom-exact-coefficient}
\end{equation}
Field-independent subtractions can be omitted inside
\([t^2]\). The factors \(2,-1,+1\) follow from
\(\beta=2[t^2]\rho_0\), and the determinant weights
\(1,-\tfrac12,+\tfrac12\).

Every coefficient on the right is defined by the
localized matrix-logarithm derivatives at zero.
Their source kernels absorb the polynomial weights,
so the sums are convergent at fixed endpoint.
For \(M=1\), \(s''\) in \(V_a(t)\) is the actual bare
Wilson Hessian on an explicit commuting constant-curvature
background. The coarse contribution is uniformly bounded
in \(M,K\) by Notebook 32. Hence the remaining ultraviolet
comparison can be made directly on the first two terms
of \eqref{cdmom-exact-coefficient}.

\subsection{What this settles and what it does not}
\label{cdmom-scope}

The formula matches an explicit constant-curvature trace
calculation to the actual Maxwell extraction, with the
actual composite average, stationary jets and measure
factors. It removes the multiplier-curvature jets on
this probe. It does not replace an arbitrary local
density by its value on an affine field: the individual
gauge identity of the chosen diagonal densities was
essential.

On a fixed periodic torus, constant nonzero Cartan flux
has quantization conditions. No analytic family of
periodic fluxes through zero is used here. Nor is an
unbounded affine field substituted into the original
bounded-field map. The identities are polynomial
extensions of localized derivatives in an infinite
joint system.

The fine vector--ghost expression still needs a bound
on its difference from a precisely normalized reference
logarithm, uniform as \(K\) grows. The fixed-endpoint
logarithm construction supplies no such bound and
computes no sign. Mixed and large-field histories,
the continuum observable construction and reconstruction
remain separate requirements for full existence.


% No preamble. Prefix mcr-. This is a continuum reference, not a lattice comparison.
\section{A completely specified massive continuum reference}
\label{mcr-section}

We compute a reference with its regulator and normalization fixed.
This supplies a target for comparison, not an identification of
the actual constrained lattice determinant with that target.
The background-field operators and heat coefficient agree with
the conventions described in \cite{vassilevichctn}, Sections 3.4,
4.2.1 and 8.3. The calculation below keeps the regulator and
the one-color Maxwell extraction explicit.

Use final physical coordinates on \(\mathbb R^4\), scalar mass
\(m>0\), and a lower proper-time cutoff \(\epsilon^2>0\).
Choose a Hermitian Cartan element \(H\) with
\(\tau(H^2)=1\). Let \(q\) run over the eigenvalues of
\(\operatorname{ad}_H\) on the complexified adjoint space,
with multiplicities. Both signs of a nonzero charge are
included. Complexification has the same trace as the
original real color space; it does not introduce another
factor of two. Put
\[
 C_H=\sum_q q^2.
\]
For one constant curvature component \(F_{12}=fH\), a charge
sector has scalar operator \(L_q=-D_\mu D_\mu\), where one
may choose \(D_2=\partial_2+iqf x_1\) and the other derivatives
ordinary. The vector reference is
\[
 (V_q)_{\mu\nu}=L_q\delta_{\mu\nu}-2iqF^{(1)}_{\mu\nu},
 \qquad F^{(1)}_{12}=f,\quad F^{(1)}_{21}=-f.
\]
The spin matrix has eigenvalues \(0,0,2qf,-2qf\).
Its sign convention has no effect on its trace.
We add \(m^2\) to both operators.

\subsection{Exact diagonal heat traces}
\label{mcr-heat-section}

Fourier transformation in \(x_2\) reduces the magnetic
plane to oscillators with eigenvalues \(|qf|(2n+1)\).
For a normalized oscillator eigenfunction, integration
of its shifted squared value gives
\[
 \int_{\mathbb R}\frac{dp_2}{2\pi}
  \left|\psi_n\!\left(x_1+\frac{p_2}{qf}\right)\right|^2
       =\frac{|qf|}{2\pi}.
\]
Thus this is a pointwise diagonal density calculation,
without a finite-box degeneracy assumption. Summing the
oscillator levels and multiplying by the heat density
of the two free coordinates gives
\begin{equation}
 k_q(s;f)=(4\pi s)^{-2}\frac{qfs}{\sinh(qfs)}.
 \label{mcr-scalar-heat}
\end{equation}
At \(qf=0\) the fraction is defined by its limit.
The vector density is
\[
 k_q(s;f)\{2+2\cosh(2qfs)\}.
\]

Define the normalized proper-time determinant correction
with weights \(-\tfrac12\log\det(V+m^2)+\log\det(L+m^2)\).
Here this phrase means the following density integral;
there is no infinite-volume determinant:
\begin{equation}
 {\cal R}_{\epsilon,m}(f)
 =\frac1{16\pi^2}\int_{\epsilon^2}^{\infty}
       \frac{ds}{s^3}e^{-m^2s}
       \sum_q\left\{\frac{qfs}{\sinh(qfs)}
                          \cosh(2qfs)-1\right\}.
 \label{mcr-reference}
\end{equation}
The sign follows from
\(\log P=-\int ds\,e^{-sP}/s\), after subtraction and
the specified proper-time cutoff. The two vector
polarizations with zero spin cancel the scalar ghost
in this combined expression.

For real \(|f|\max_q|q|<m^2\), the massive vector has
positive bottom and the integral converges at infinity.
On a smaller real interval, differentiation twice at
zero is justified by a common exponentially decaying
majorant. In particular, the notation \([f^2]\) below
denotes half the second real derivative at zero.
We do not assert a convergent complex Taylor series
in a nonzero constant magnetic field.

\begin{proposition}[Normalized reference coefficient]
\label{mcr-coefficient}
Set \(E_1(x)=\int_x^\infty e^{-u}\,du/u\). In the
normalization \(\beta=2[f^2]\) used in this notebook,
\begin{equation}
 \beta_{\rm ref}(\epsilon,m)
       =\frac{11C_H}{48\pi^2}E_1(m^2\epsilon^2).
 \label{mcr-beta}
\end{equation}
In particular, with \(m=1\), \(\epsilon=K^{-1}\),
\begin{equation}
 \left|\beta_{\rm ref}(K^{-1},1)
           -\frac{11C_H}{24\pi^2}\log K\right|
 \le \frac{11C_H}{48\pi^2}(1+e^{-1}),\qquad K\ge1.
 \label{mcr-finite-part}
\end{equation}
\end{proposition}

\begin{proof}
The scalar magnetic factor and the spin factor give
\[
 \frac{x}{\sinh x}\cosh(2x)
       =1+\frac{11}{6}x^2+O(x^4).
\]
Differentiating \eqref{mcr-reference} twice at zero
therefore gives
\[
 [f^2]{\cal R}_{\epsilon,m}
    =\frac{11C_H}{96\pi^2}
           \int_{\epsilon^2}^{\infty}\frac{e^{-m^2s}}s\,ds.
\]
The substitution \(u=m^2s\) proves \eqref{mcr-beta}.
For \(0<x\le1\),
\[
 E_1(x)-\log(1/x)
   =\int_x^1\frac{e^{-u}-1}{u}\,du
                  +\int_1^\infty\frac{e^{-u}}u\,du.
\]
The absolute value is at most \(1+e^{-1}\). This proves
the asserted uniform finite part.
\end{proof}

For the repository's \(SU(N)\) normalization
\(\tau=\operatorname{Tr}_{\rm fund}/N\), write
\(H=\operatorname{diag}(h_1,\ldots,h_N)\). Then
\(\sum_i h_i=0\), \(\sum_i h_i^2=N\), and the
nonzero-color charges are \(h_i-h_j\), \(i\ne j\);
the \(N-1\) Cartan charges vanish. Consequently
\[
 C_H=\sum_{i,j}(h_i-h_j)^2
              =2N\sum_i h_i^2=2N^2.
\]
The coefficient of \(\log K\) is thus
\(11N^2/(12\pi^2)\), or \(11/(3\pi^2)\) for \(SU(2)\),
in this one-normalized-color convention.
This is not a sum of \(N^2-1\) different Maxwell extractions.

\subsection{Scale and regulator bookkeeping}
\label{mcr-scaling-section}

The same continuum reference can be written in current
fine units, with cutoff \(s\ge1\), mass \(K^{-1}\), and
curvature \(f/K^2\). Multiplication of that density by
the block volume \(K^4\), followed by \(s=K^2u\), gives
exactly \eqref{mcr-reference} with \(\epsilon=K^{-1}\),
\(m=1\). Thus there is no additional factor of \(K^4\)
in its extracted coefficient. This is a change of units
within the continuum reference, not a replacement of a
lattice heat kernel by a continuum kernel.

Fixed positive changes of these continuum regulators
have a bounded effect. For \(c,m>0\),
\begin{equation}
 |\beta_{\rm ref}(c/K,m)-\beta_{\rm ref}(1/K,1)|
 \le\frac{11C_H}{24\pi^2}|\log(cm)|.
 \label{mcr-regulator-change}
\end{equation}
Indeed, the difference of the two \(E_1\) integrals is
bounded by the integral of \(du/u\) between their lower
limits. Uniformity follows when \(cm\) remains in a
fixed compact subinterval of \((0,\infty)\).
This estimate does not compare a based ghost to an
unbased massive ghost, change a lattice constraint, or
control a massless zero-mode limit.



% No preamble. Prefix ahc-. Actual fixed-endpoint identity; uniform estimates conditional.
\section{The actual heat integral and a sufficient comparison estimate}
\label{ahc-section}

The Cartan identity also puts the missing comparison into
heat-kernel form without changing the lattice operators.
Use \(G(t),V_a(t)\) from \eqref{cdmom-probe-matrices}.
Let \(\operatorname{Tr}_{0,S}\) be the diagonal color trace
over noncenter scalar sites in block zero and
\(\operatorname{Tr}_{0,E}\) that over current bonds with
initial site in block zero. Define the formal coefficient
\begin{equation}
 h_{M,K}(s)=[t^2]\left\{
   \operatorname{Tr}_{0,E}e^{-sV_a(t)}
                  -2\operatorname{Tr}_{0,S}e^{-sG(t)}
                         \right\}.
 \label{ahc-h-definition}
\end{equation}
Every derivative here is evaluated at the flat positive
operators and on the polynomial stationary jets. No
heat evolution at a nonzero unbounded background is
used. Put
\[
 b_c(M,K)=[t^2]\sum_{c_-=0}
           \operatorname{tr}_{\mathfrak g}(\log P(t))_{cc}.
\]
Notebook 32 bounds \(|b_c|\) uniformly in \(M,K\).

\begin{proposition}[Exact actual heat representation]
\label{ahc-identity}
In each joint infinite system at fixed \(K\),
\begin{equation}
 \beta(D_{M,K})=b_c(M,K)
                   +\int_0^\infty h_{M,K}(s)\,\frac{ds}{s}.
 \label{ahc-exact}
\end{equation}
The integral is absolutely convergent for fixed \(K\).
The bounds establishing this statement may depend on \(K\).
\end{proposition}

\begin{proof}
First use bounded compactly supported output directions.
For a positive flat matrix family \(P(t)\), the local
functional calculus gives
\[
 [t^2](\log P(t))_{ii}
    =-\int_0^\infty [t^2](e^{-sP(t)})_{ii}\,\frac{ds}{s}.
\]
One can obtain this by differentiating the difference
of the proper-time formulas for \(P(t)\) and \(P(0)\);
their scalar reference term cancels.
To justify the same formula on polynomial jets, write
\(P(t)=P_0+tP_1+t^2P_2+O(t^3)\). Its heat coefficient is
\begin{align}
 [t^2]e^{-sP(t)}
 ={}&-\int_0^s e^{-(s-u)P_0}P_2e^{-uP_0}\,du \notag\\
 &+\int_{0<u<v<s}
 e^{-(s-v)P_0}P_1e^{-(v-u)P_0}P_1e^{-uP_0}\,du\,dv .
 \label{ahc-duhamel}
\end{align}
The \(P_j\) kernels are polynomially growing at worst,
with exponential localization in their relative
coordinates. This follows by inserting the affine
source jets in the established rooted derivative kernels.

The fixed-\(K\) flat gap and exponential matrix bound
give exponentially weighted semigroup estimates.
Indeed bounded truncated distance conjugation changes
\(P_0\) in operator norm by less than half its gap if
the exponent is small enough. The variation-of-constants
bound then gives
\(\|e^{-sP_0}\|_{\rm weighted}\le e^{-c_Ks}\).
Spending smaller spatial exponents absorbs the
polynomial factors in \eqref{ahc-duhamel}, after
anchoring at the diagonal block. Its trace is bounded by
\[
 C_K(s+s^2)e^{-c'_Ks},\qquad s>0,
\]
with constants uniform in volume and, where the
upstream bounds are, preceding depth. For small \(s\)
the same estimate follows directly from the one and
two insertion formulas. These bounds justify the
proper-time integral, polynomial-jet extension and
local infinite-volume passage by domination.
Substituting the two log coefficients in
\eqref{cdmom-exact-coefficient} proves
\eqref{ahc-exact}, with signs \(+1,-2\) in
\eqref{ahc-h-definition}.
\end{proof}

\subsection{A uniform short-time bound for the actual based ghost}
\label{ahc-ghost-short-section}

One part of the short-time comparison can be proved directly.
It uses the actual based scalar domain and has constants
independent of the preceding depth and endpoint.

\begin{proposition}[Actual ghost ultraviolet bound]
\label{ahc-ghost-short}
In four dimensions put
\[
 h^g_K(s)=-2[t^2]\operatorname{Tr}_{0,S}e^{-sG(t)}.
\]
For every permitted \(K\), on the affine probe already
identified with the actual stationary jets,
\begin{equation}
 0\le h^g_K(s)\le C_H\{(8s)^4+6(8s)^3+7(8s)^2+8s\},
                     \qquad s>0.
 \label{ahc-ghost-poisson}
\end{equation}
Consequently
\begin{equation}
 \int_0^1 |h^g_K(s)|\,\frac{ds}{s}\le2280C_H.
 \label{ahc-ghost-uv}
\end{equation}
Equivalently, the based-ghost part of the rescaled
short-time integral over \(0<u<K^{-2}\) is bounded
by this same constant.
\end{proposition}

\begin{proof}
On the complexified charge-\(q\) space, the actual based
operator is the covariant nearest-neighbor Laplacian
with Dirichlet values at every \(K\)-center. Hence
\[
 G_q(t)=8I-A_q(t),
\]
where \(A_q\) is the covariant adjacency operator on
noncenter sites. The diagonal is still \(8\), also
at sites adjacent to a removed center: it comes from
the full link norm of \(D_b\), with the center value
set to zero.

The diagonal of \(A_q(t)^n\) is a sum over length-\(n\)
closed nearest-neighbor walks avoiding the centers.
On the commuting affine field each walk has phase
\[
       \exp\{iqt\,{\cal A}_{\mu\nu}(\gamma)/K^2\},
\]
where \({\cal A}_{\mu\nu}(\gamma)\) is its oriented
algebraic area in the selected coordinate plane.
The contribution from a constant potential cancels
on a closed walk. Taking the origin of the area
coordinate at the starting site gives the elementary
bound \(|{\cal A}_{\mu\nu}(\gamma)|\le n^2\).
Thus the block trace coefficient is exactly
\begin{equation}
 h^g_K(s)=\frac{C_H}{K^4}e^{-8s}
   \sum_{\substack{z\ {\rm noncenter}\\z\ {\rm in\ block}\ 0}}
   \sum_{n\ge0}\frac{s^n}{n!}
   \sum_{\substack{\gamma:z\to z,\ |\gamma|=n\\
                    \gamma\ {\rm avoids\ centers}}}
                         {\cal A}_{\mu\nu}(\gamma)^2.
 \label{ahc-ghost-walk}
\end{equation}
All sums converge absolutely at each fixed \(s\);
the factorial absorbs the factor \(n^4\). This
also justifies extracting the formal quadratic
coefficient without imposing periodic magnetic flux.
Every term is nonnegative.

There are at most \(K^4\) starting sites and at most
\(8^n\) walks from each. It follows that
\[
 h^g_K(s)\le C_H e^{-8s}
                   \sum_{n\ge0}\frac{(8s)^n n^4}{n!}.
\]
The fourth moment of this exponential generating
series is \(x^4+6x^3+7x^2+x\) after multiplication
by \(e^{-x}\); it follows by applying
\((x\,d/dx)^4\) to \(e^x\). This proves
\eqref{ahc-ghost-poisson}. Integrating its four
monomials from zero to one with measure \(ds/s\)
gives
\[
 C_H\left(\frac{8^4}{4}
             +\frac{6\,8^3}{3}+\frac{7\,8^2}{2}+8\right)
                         =2280C_H.
\]
\end{proof}

This bound is uniform in \(M\) because the actual
stationary ghost jets on this probe have already
been identified with the same affine covariant
Laplacian for every \(S_M\). No estimate for a
separate massive inverse is used. The polynomial
upper bound in \eqref{ahc-ghost-poisson} is not a
useful large-time bound and does not settle the
long-time comparison.

\subsection{A quantitative target in physical heat time}
\label{ahc-target-section}

Set \(J_{M,K}(u)=h_{M,K}(K^2u)\), and
\(\alpha_H=11C_H/(48\pi^2)\).
The continuum calculation above has reference heat
coefficient
\[
 J_K^{\rm ref}(u)=\alpha_H e^{-u}
                         {\bf1}_{\{u\ge K^{-2}\}}.
\]
Thus the following is an exact difference identity,
not yet a bound:
\begin{equation}
 \beta(D_{M,K})-\beta_{\rm ref}(K^{-1},1)
 =b_c(M,K)+\int_0^\infty
          \{J_{M,K}(u)-J_K^{\rm ref}(u)\}\,\frac{du}{u}.
 \label{ahc-difference}
\end{equation}
It fixes the actual operator, regulator, sign, density,
and rescaling in a proposed comparison.

\begin{lemma}[Sufficient bounds on three heat-time intervals]
\label{ahc-band-lemma}
Suppose, for all required \(M,K\) and structural systems,
there are constants \(C,\delta>0\), independent of them,
such that
\begin{align}
 &\int_0^{K^{-2}}|J_{M,K}(u)|\,\frac{du}{u}\le C,
                                      \label{ahc-uv-bound}\\
 &|J_{M,K}(u)-\alpha_H|
       \le C\{(K^2u)^{-\delta}+u^\delta\},
          \quad K^{-2}\le u\le1,       \label{ahc-mid-bound}\\
 &\int_1^\infty|J_{M,K}(u)|\,\frac{du}{u}\le C.
                                      \label{ahc-ir-bound}
\end{align}
Then, with the already uniform coarse bound \(C_c\),
\begin{equation}
 \left|\beta(D_{M,K})
       -\frac{11C_H}{24\pi^2}\log K\right|
           \le C_c+2C+\frac{2C}{\delta}.
 \label{ahc-log-conclusion}
\end{equation}
\end{lemma}

\begin{proof}
Integrate the exact identity \eqref{ahc-exact} after
the time rescaling. The constant \(\alpha_H\) on the
middle interval integrates to \(2\alpha_H\log K\).
Each error in \eqref{ahc-mid-bound} integrates to
at most \(C/\delta\). The other intervals and the
coarse term give the remaining terms.
\end{proof}

The full bounds \eqref{ahc-uv-bound}--\eqref{ahc-ir-bound}
are not asserted here for the actual family uniformly in
\(K\). Proposition~\ref{ahc-ghost-short} supplies the
ghost part of the first bound; its vector part remains.
The fixed-endpoint heat representation by itself gives
only constants depending on \(K\). The middle interval asks for a
continuum coefficient on heat lengths between the
fine spacing and the final block size. The last
interval is where the soft pinned-ghost modes are
particularly relevant: the scalar gap alone may be
of order \(K^{-4}\) in fine units, so rescaling by
\(K^2\) does not give uniform damping. The combined
coefficient must be estimated, or a weaker signed
integral comparison must be proved. Bounding the
separate inverse operators by a physical mass is
not an available shortcut.

For \(M>1\), \(V_a\) contains the actual \(S_M\)
Hessian. Uniformity over that family is part of the
conditions, not a consequence of computing the bare
Wilson reference. A bounded comparison only at \(M=1\)
would not by itself give the varying-depth margin
needed by the finite drift-transfer theorem.



\section{The remaining scale-dependent comparison}
\label{ctn-next}

For the bare composite, the new identification reduces
the leading coefficient, up to the uniformly bounded
coarse effective-Hessian term, to the constant-curvature
jets of
\[
 \log\det(D_b^TD_b)
 -\frac12\log\det
       \bigl(S_{\mathrm W}''+D_bD_b^T+aC_K^TC_K\bigr).
\]
This is a statement about the actual diagonal trace density
and its Maxwell moment. The formal Cartan probe makes
the constrained-current curvature vanish in the required
jets. It does not make either determinant individually
uniformly massive at the final scale.

The next estimate must compare this combined expression
to a fully specified background-Feynman vector and ghost
reference with infrared scale \(K^{-1}\), preserving its
coarse-coordinate, color and action normalization. In
particular, changing the based scalar domain, the vector
constraint completion and the infrared regulator must
be controlled together. The based-ghost Rayleigh bound
from Notebook 32 still rules out a separate inverse
bound of order \(K^2\).

A sufficient result would give a constant \(C\), independent
of \(K\), bounding the difference of the extracted
Maxwell coefficients. The specified continuum reference
now has an explicit positive logarithmic law and bounded
finite part. The exact actual heat representation identifies
the remaining comparison integral; the three interval
estimates in Lemma~\ref{ahc-band-lemma} would bound it.
Those estimates do not follow from the fixed-\(K\)
convergence used here. The actual based-ghost part of
the short-time interval is now bounded uniformly by
the closed-walk argument; its vector counterpart and
the other intervals remain. Uniformity over
the structural systems needed by the finite drift-transfer
theorem must also be retained.

Once positive averaged drift is established with that
scope, the conditional finite shooting theorem applies.
The full four-dimensional existence objective still
requires the other field sectors, the exact global
measure representation, sourced continuum bounds,
nontriviality, Euclidean symmetry, positivity and
reconstruction.

\begin{thebibliography}{9}
\bibitem{vassilevichctn}
D.~V.~Vassilevich,
\emph{Heat kernel expansion: user's manual},
Physics Reports 388 (2003), 279--360,
\href{https://arxiv.org/abs/hep-th/0306138}{arXiv:hep-th/0306138}.
Sections 3.4, 4.2.1 and 8.3 specify the background vector
and ghost operators, their combined heat coefficient,
and the constant magnetic scalar heat kernel.
This source is not used as a proof of the actual constrained
lattice comparison.
\bibitem{priorctn}
\emph{Toward a four-dimensional Yang--Mills construction},
research Notebooks 15--32 in this workspace.
The actual local minimum, strict mixed kernels, classical
envelope, constrained Gaussian, gauge identities, determinant
bridge and Maxwell normalization are the inputs used here.
\bibitem{sourcectn}
E.~Sussman and collaborators,
\href{https://github.com/ethanwsussman/Yang--Mills-LLM-experiment}
{Yang--Mills-LLM-experiment}, pinned deposited accounts at
commit \texttt{e392ef45013a581664bab76b90867170b30fe10e}.
The centered ordinary average and its selected
endpoint-equivariant mean are kept unchanged.
\end{thebibliography}
\end{document}
